\documentclass[11pt]{article}
\usepackage[utf8]{inputenc}
\usepackage[T1]{fontenc}
\usepackage[a4paper,margin=1.8cm]{geometry}
\usepackage{booktabs}
\usepackage{longtable}
\usepackage{array}
\usepackage{xcolor}
\usepackage{hyperref}
\usepackage{pdflscape}

\hypersetup{colorlinks=true, linkcolor=blue, urlcolor=blue}
\setlength{\parindent}{0pt}
\renewcommand{\arraystretch}{1.18}

\newcolumntype{L}[1]{>{\raggedright\arraybackslash}p{#1}}
\newcommand{\code}[1]{\texttt{\small #1}}

\begin{document}

\section*{Findings and Evidence: Platform Membership as Marketing-Mix Moderator}

This table summarizes the interpretation-relevant findings from the clean
analysis module. The main outcome is transformed capacity utilisation
(\code{ur\_boxcox}). The main fixed-effect strategies are \code{market\_fe}
(city, date, and operator fixed effects) and \code{unit\_fe}
(city-operator and date fixed effects). Standard errors are clustered by
\code{city\_operator}. The referenced files are all included in this repository.

\begin{landscape}
\small
\begin{longtable}{L{0.22\linewidth} L{0.32\linewidth} L{0.38\linewidth}}
\caption{Relevant findings, empirical reference values, and interpretation.}\\
\toprule
\textbf{Finding} & \textbf{Reference value} & \textbf{Interpretation sentence} \\
\midrule
\endfirsthead
\toprule
\textbf{Finding} & \textbf{Reference value} & \textbf{Interpretation sentence} \\
\midrule
\endhead
\bottomrule
\endfoot

No robust universal platform bonus &
\code{platform\_any}: market-FE coef $=0.0385$, $p=0.0468$; unit-FE coef $=0.0040$, $p=0.788$ &
Platform membership should not be framed as an unconditional utilisation booster because the positive market-FE association disappears once city-operator fixed effects are used. \\

Platform type is empirically heterogeneous &
Mean \code{ur}: none $=0.00755$, FreeNow-only $=0.00845$, local MaaS-only $=0.01775$, multi-platform $=0.00816$, weak exposure-only $=0.00691$ &
The data support treating platform membership as a set of distinct platform contexts rather than a homogeneous binary state. \\

Local MaaS has the highest descriptive utilisation &
Local MaaS-only: 1,686 rows, 18 city-operators, mean \code{ur} $=0.01775$ versus none $=0.00755$ &
Local MaaS integrations appear descriptively associated with substantially higher capacity utilisation, but causal interpretation still requires the moderation models. \\

Weak exposure is not equivalent to membership &
Weak exposure-only: 305 rows, 2 city-operators, mean \code{ur} $=0.00691$; unit-FE main term coef $=-0.0617$, $p<0.001$ &
Weak exposure should be separated from real bookable/payable platform membership because it behaves empirically like a different condition. \\

Core moderation: access pricing friction &
\code{unlock\_fee $\times$ any\_platform}: market-FE coef $=-0.0715$, $p<0.001$; 9 negative significant estimates across 16 model variants &
The strongest story is that platform membership changes marketing effectiveness by making access-pricing frictions more harmful for capacity utilisation. \\

Local MaaS amplifies unlock-fee friction &
\code{unlock\_fee $\times$ local\_maas\_only}: market-FE coef $=-0.1075$, $p<0.001$; 9 negative significant estimates across 16 variants &
The pricing-friction result is especially pronounced in local MaaS contexts, suggesting that integrated mobility platforms may heighten sensitivity to fixed access costs. \\

Unlock fee is more central than per-minute price &
\code{unlock\_fee $\times$ any\_platform}: coef $=-0.0715$, $p<0.001$; \code{price/min $\times$ any\_platform}: coef $=-0.0049$, $p=0.818$ in the same market-FE model &
The relevant pricing mechanism is not the variable usage price but the fixed access friction that users face before a trip starts. \\

FreeNow-only is not a robust positive performance context &
\code{large\_aggregator\_only}: market-FE main coef $=-0.0461$, $p=0.0559$; unit-FE main coef $=0.0006$, $p=0.970$ &
FreeNow-only membership does not provide a stable direct capacity-utilisation advantage once richer fixed effects are applied. \\

FreeNow promotion effect is suggestive but not robust &
\code{promotion $\times$ FreeNow-only}: market-FE coef $=0.1304$, $p<0.001$; unit-FE coef $=0.0002$, $p=0.987$ &
There is suggestive evidence that promotions work better in the FreeNow context, but it should be interpreted cautiously because it is not robust to city-operator fixed effects. \\

Coverage effects are not universally platform-enhanced &
\code{coverage $\times$ any\_platform}: market-FE coef $=-0.0245$, $p=0.112$; unit-FE coef $=0.0014$, $p=0.862$ &
Platform membership does not generally make simple service coverage more effective for capacity utilisation. \\

Multi-platform coverage has a unit-FE signal &
\code{coverage $\times$ multi\_platform}: unit-FE coef $=0.0516$, $p=0.006$; market-FE coef $=-0.0045$, $p=0.840$ &
Coverage may matter more in multi-platform contexts within the same city-operator unit, but this result is narrower than the unlock-fee finding. \\

Exclusive coverage can turn negative in platform contexts &
\code{exclusive\_coverage $\times$ FreeNow-only}: market-FE coef $=-0.0483$, $p=0.0065$; local MaaS unit-FE coef $=-0.0731$, $p<0.001$; multi-platform unit-FE coef $=-0.0664$, $p<0.001$ &
Exclusive territorial coverage is not consistently beneficial under platform membership and may reflect lower-comparison or lower-demand contexts. \\

Competition mechanism is plausible but exploratory &
\code{unlock\_fee $\times$ any\_platform $\times$ competition}: market-FE coef $=-0.0330$, $p=0.099$; no-weak-exposure market-FE coef $=-0.0394$, $p=0.060$ &
Competition provides a plausible mechanism for the pricing-friction story, but the evidence is not stable enough to serve as the main claim. \\

DML partially supports the mechanism but is not decisive &
HGB-DML no-weak-exposure: \code{unlock\_fee $\times$ local\_only $\times$ competition} coef $=-0.0353$, $p=0.0057$; RF-DML signs are not fully consistent &
Double Machine Learning offers some support for competition-amplified pricing friction, but estimator sensitivity means the DML evidence should be framed as robustness probing rather than the core result. \\

Capacity utilisation is the appropriate performance lens &
Primary outcome is \code{ur\_boxcox}; robustness outcome is \code{log\_trip\_count\_day}; all main tables report 14,514 all-sample observations and 97 city-operators &
The contribution is about provider-side efficiency of deployed capacity, not only about demand volume or trip counts. \\

\end{longtable}
\end{landscape}

\section*{Naive interpretation}

A naive reading of platform membership would suggest that joining a platform
should improve provider performance by expanding reach, increasing visibility,
and giving providers access to a larger pool of potential users. Our results do
not support this simple demand-expansion interpretation. Instead, the evidence
suggests that platform membership changes the conditions under which marketing
instruments translate into capacity utilisation. In particular, platform contexts
appear to make fixed access-pricing frictions more consequential: unlock fees are
more negatively associated with capacity utilisation when providers are embedded
in platform-mediated markets, and this pattern is especially visible for local
MaaS integrations. This implies that platforms do not merely add demand; they
also reshape the choice environment in which users compare providers and evaluate
the cost of starting a trip.

The paper-worthy interpretation is therefore that platform membership alters the
effectiveness of the marketing mix rather than producing a universal performance
premium. Providers cannot treat platforms as neutral reach channels. Once a
provider becomes visible or bookable through a platform, access costs such as
unlock fees become more salient and may reduce the efficiency with which deployed
capacity is used. Local MaaS platforms appear to be the clearest setting for this
mechanism, while FreeNow-only membership does not show a robust direct
utilisation advantage. The contribution is thus to shift the provider-side
platform discussion from whether platform membership pays off on average to when
and through which marketing instruments it changes provider performance.

\end{document}
